\chapter{线段上的积分}\label{ch:b1:integration}

高中分册中的积分建立在直观取得的面积之上。
本章从头构造积分：先对阶梯函数给出构造，
此时积分为有限和；再对连续（以及分段连续）函数，
借助一致逼近来构造 ——
这正是海涅定理（\cref{thm:b1:continuity:heine}）大显身手之处。
随后，微积分基本定理把这一构造与原函数联系起来，
黎曼和则把它与离散平均
联系起来。

全文约定 $a < b$ 为实数。

\section{阶梯函数}

\begin{definition}\label{def:b1:integration:step}
映射 $\varphi \colon \intcc{a}{b} \to \R$ 称为\emph{阶梯函数}\index{阶梯函数}，
当存在剖分 $a = x_0
< x_1 < \dots < x_n = b$，使 $\varphi$ 在每个开区间上恒等于
$c_i$，开区间为 $\intoo{x_{i-1}}{x_i}$，
节点处的取值不作要求。其积分定义为
\[
\int_a^b \varphi = \sum_{i=1}^{n} c_i\,(x_i - x_{i-1}),
\]
该积分与所取剖分无关（将两个剖分取公共加细，
两边在加细下均保持不变）。
\end{definition}

\begin{proposition}\label{prop:b1:integration:stepprops}
对阶梯函数，积分是线性的，且保序（$\varphi \leq
\psi \implies \int\varphi \leq \int\psi$），并满足夏勒关系
$\int_a^b = \int_a^c + \int_c^b$（$a < c < b$）。
\end{proposition}

\begin{proof}
证明的关键在于\emph{加细不变性}，即定义中已述的性质：
在剖分中插入一个额外节点 $t \in \intoo{x_{i-1}}{x_i}$，
则将项 $c_i(x_i - x_{i-1})$ 替换为
$c_i(t - x_{i-1}) + c_i(x_i - t)$ --- 仍是同一个数 ---，所以
积分不因任何有限加细而改变。现取
$\varphi$ 具有剖分 $\sigma$，$\psi$ 具有剖分 $\sigma'$，
在公共加细 $\sigma \cup \sigma'$ 上，两者
都是具有\emph{相同}节点的阶梯函数，
在每个小区间上 $\varphi + \lambda\psi$ 恒等于 $c_i + \lambda
d_i$，线性性归结为有限和的线性性。保序性：
在每个小区间上 $c_i \leq d_i$，故 $\sum c_i \Delta_i \leq \sum
d_i \Delta_i$（长度 $\Delta_i \geq 0$）。夏勒关系：插入
节点 $c$，并在该处拆分和式。
\end{proof}

\section{连续函数的积分}

\begin{theorem}[一致逼近]\label{thm:b1:integration:approx}
设 $f$ 在 $\intcc{a}{b}$ 上连续。对任意 $\varepsilon > 0$，
存在阶梯函数 $\varphi, \psi$，使得
\[
\varphi \leq f \leq \psi
\qquad\text{且}\qquad
\psi - \varphi \leq \varepsilon \text{ 在 } \intcc{a}{b}.
\]
\end{theorem}

\begin{proof}
由海涅定理（\cref{thm:b1:continuity:heine}），$f$ 一致连续：
取 $\delta$（与 $\varepsilon$ 相应），并取细度 $< \delta$ 的剖分
（例如等距剖分，取 $n > \frac{b -
a}{\delta}$ 个小区间）。在每个闭子区间 $\intcc{x_{i-1}}{x_i}$ 上，$f$
取得最小值 $m_i$ 与最大值 $M_i$
（见\cref{thm:b1:continuity:evt}），且 $M_i - m_i \leq \varepsilon$
（两个极值点的距离不超过 $\delta$）。定义 $\varphi = m_i$
及 $\psi = M_i$，在 $\intoo{x_{i-1}}{x_i}$ 上
（在节点处 $\varphi = \psi = f$）。
\end{proof}

\begin{theorem}[积分的定义]\label{thm:b1:integration:def}
设 $f$ 在 $\intcc{a}{b}$ 上连续。两个数
\[
I_-(f) = \sup\Bigl\{\int_a^b \varphi : \varphi \text{ 为阶梯函数},\
\varphi \leq f\Bigr\},
\qquad
I_+(f) = \inf\Bigl\{\int_a^b \psi : \psi \text{ 为阶梯函数},\ \psi \geq
f\Bigr\}
\]
二者相等；其公共值即为
\emph{积分}\index{积分} $\int_a^b f$（亦记作 $\int_a^b
f(t)\,\dd t$）。它在阶梯函数上与前面的定义一致，
并将其推广到分段连续函数，
方法是在间断点处切分 $\intcc{a}{b}$（以夏勒关系
作为此处的定义）。
\end{theorem}

\begin{proof}
两个集合均非空（$f$ 有界），且每个下阶梯积分
都 $\leq$ 每个上阶梯积分（阶梯函数上的保序性）。于是
$I_-(f) \leq I_+(f)$。由\cref{thm:b1:integration:approx}，对每个
$\varepsilon$，存在一对阶梯函数，使 $\int\psi - \int\varphi \leq
\varepsilon(b - a)$：上确界与下确界被挤压到一起，$I_- =
I_+$。
\end{proof}

\begin{example}[分段连续，并不复杂]\label{ex:b1:integration:piecewise}
在 $\intcc{0}{3}$ 上，取整函数其实是一个阶梯函数，
只是形式上稍加伪装：在其跳跃点处切分，
\[
\int_0^3 \lfloor t \rfloor\,\dd t
= \int_0^1 0 + \int_1^2 1 + \int_2^3 2 = 0 + 1 + 2 = 3 ,
\]
而\emph{在}跳跃点 $1, 2$ 处的取值无关紧要：
改变函数在有限多个点上的取值，不改变任何积分，
（用于夹逼的阶梯函数不受影响）。这正是
“分段连续”推广的全部内容：在有限个间断点处切开，
对每个连续片段分别积分，
再相加 --- 以夏勒关系为定义。
\end{example}

\begin{example}[定义只需计算一次]\label{ex:b1:integration:defcompute}
设 $f(x) = x$ 在 $\intcc{0}{1}$ 上，并将其等分为 $n$ 段。
在各段上取常值的最优阶梯函数分别为 $\varphi =
\frac{k-1}{n}$ 与 $\psi = \frac kn$，在第 $k$ 段上，并有
\[
\int_0^1 \varphi = \sum_{k=1}^{n} \frac{k-1}{n}\cdot\frac1n
= \frac{n-1}{2n},
\qquad
\int_0^1 \psi = \sum_{k=1}^{n} \frac{k}{n}\cdot\frac1n
= \frac{n+1}{2n} .
\]
每一个下积分都 $\leq I_-(f) \leq I_+(f) \leq$
每一个上积分，于是 $\frac{n-1}{2n} \leq I_-(f) \leq I_+(f) \leq
\frac{n+1}{2n}$ 对一切 $n$ 成立，两者都夹向 $\frac12$，
而 $\int_0^1 x\,\dd x = \frac12$ 直接由定义可得。
小结：这是第一次、也是最后一次
按定义计算积分 --- 下文的微积分基本定理
以一次原函数查找取代一切此类计算，这正是
本章的核心意义。
\end{example}

\begin{theorem}[性质]\label{thm:b1:integration:props}
设 $f, g$ 在 $\intcc{a}{b}$ 上连续（或分段连续），且
$\lambda \in \R$：
\begin{enumerate}
  \item 线性：$\int (f + \lambda g) = \int f + \lambda \int g$；
  \item 单调性：$f \leq g \implies \int_a^b f \leq \int_a^b g$；
        且 $\bigl|\int_a^b f\bigr| \leq \int_a^b \abs f \leq (b -
        a)\, \sup\abs f$；
  \item 沙勒关系：$\int_a^b = \int_a^c + \int_c^b$（
        约定 $\int_b^a = -\int_a^b$，对积分
        上下限的任意次序均成立）；
  \item 严格正性：若 $f$ 连续，$f \geq 0$，且
        $\int_a^b f = 0$，则 $f = 0$ 在
        $\intcc{a}{b}$ 上处处成立。
\end{enumerate}
\end{theorem}

\begin{proof}
(1)--(3) 由阶梯函数出发，借助
上、下确界的定义过渡到极限。线性性需细究一次：给定
$\varepsilon > 0$，取 $\varphi_f \leq f \leq \psi_f$ 与
$\varphi_g \leq g \leq \psi_g$，且夹逼之差 $\leq \varepsilon$
（见\cref{thm:b1:integration:approx}）。对 $\lambda \geq 0$，
$\varphi_f + \lambda\varphi_g \leq f + \lambda g \leq \psi_f +
\lambda\psi_g$ 构成阶梯函数的夹逼，其差 $\leq
(1 + \lambda)\varepsilon$，其阶梯积分等于 $\int
\varphi_f + \lambda\int\varphi_g$ 等
（见\cref{prop:b1:integration:stepprops}）：令 $\varepsilon \to
0$ 时，$\int(f + \lambda g)$ 被夹逼至 $\int f + \lambda\int g$。
对 $\lambda < 0$，乘以 $\lambda$ 会\emph{颠倒}
$g$ 的夹逼 --- $\lambda g$ 的下阶梯函数是
$\lambda \psi_g$ --- 同样的夹逼过程照常进行，角色
对调。不等式 $\abs{\int f} \leq \int \abs f$ 由
$-\abs f \leq f \leq \abs f$ 及单调性得出。

(4) 用逆否命题证明：若 $f(x_0) = m > 0$，由连续性，存在
长为 $\eta > 0$ 的子区间，在其上 $f \geq \frac m2$；
在该子区间上取 $\frac m2$、其余处取 $0$ 的阶梯函数 $\leq f$，
于是 $\int f \geq \frac{m\eta}{2} > 0$。
\end{proof}

\begin{example}[沙勒关系的应用：含绝对值的积分]\label{ex:b1:integration:chasleswork}
求带绝对值的积分时，在符号改变处分段。
\[
\int_0^2 \abs{x - 1}\,\dd x
= \int_0^1 (1 - x)\,\dd x + \int_1^2 (x - 1)\,\dd x
= \frac12 + \frac12 = 1 ,
\]
同样，对 $\intcc{0}{2\pi}$，在 $\pi$ 处切开：
\[
\int_0^{2\pi} \abs{\sin t}\,\dd t
= \int_0^{\pi} \sin t\,\dd t - \int_{\pi}^{2\pi} \sin t\,\dd t
= 2 + 2 = 4 ,
\]
而 $\int_0^{2\pi} \sin t\,\dd t = 0$，抵消确实存在，
这正是为何严格正性的结论
（见\cref{thm:b1:integration:props}~(4)）附有假设 $f
\geq 0$ --- 若无此假设，积分为零
便无法说明关于 $f$ 的任何事情。小结：$\int \abs f$ 度量
\emph{面积}，$\int f$ 度量\emph{有符号的代数和}；
不等式 $\abs{\int f} \leq \int\abs f$ 精确记录了
抵消所能抹去的东西。
\end{example}

\section{微积分基本定理}

\begin{theorem}[微积分基本定理]\label{thm:b1:integration:ftc}
设 $f$ 在区间 $I$ 上连续，且 $a \in I$。函数
\[
F(x) = \int_a^x f(t)\, \dd t
\]
为 $C^1$ 类，定义于 $I$ 上，且 $F' = f$：每个连续函数
在区间上都有原函数。由此，对 \emph{任一}
$G$ 为 $f$ 的原函数：\index{微积分基本定理}
\[
\int_a^b f(t)\,\dd t = G(b) - G(a) .
\]
\end{theorem}

\begin{proof}
取定 $x_0 \in I$ 与 $\varepsilon > 0$，由 $x_0$ 处的连续性，存在
$\delta$，使得 $\abs{f(t) - f(x_0)} \leq \varepsilon$ 当 $\abs{t -
x_0} \leq \delta$ 时成立。对于 $0 < \abs{h} \leq \delta$（且 $x_0 + h \in
I$），由沙勒关系得
\[
\frac{F(x_0 + h) - F(x_0)}{h} - f(x_0)
= \frac 1h \int_{x_0}^{x_0+h} \bigl(f(t) - f(x_0)\bigr)\dd t ,
\]
其绝对值 $\leq \frac{1}{\abs h}\cdot \abs h\,
\varepsilon = \varepsilon$（由性质 (2)，对积分
上下限的两种次序均成立）。故 $F'(x_0) = f(x_0)$；$F' = f$ 连续，$F$ 为
$C^1$ 类。若 $G' = f$ 亦成立，则在区间上 $(G - F)' = 0$，于是 $G =
F + c$（见\cref{cor:b1:derivative:monotone}），且 $G(b) - G(a) = F(b)
- F(a) = \int_a^b f$。
\end{proof}

\begin{example}[计算前先看对称性]\label{ex:b1:integration:parity}
在对称区间上，奇偶性便可派上用场：若 $f$ 为奇函数，
作代换 $t \mapsto -t$，则 $\int_{-a}^{0} f$ 化为
$-\int_0^a f$，于是
\[
\int_{-a}^{a} f(t)\,\dd t = 0 ;
\qquad\text{若 $f$ 为偶函数，}\quad
\int_{-a}^{a} f = 2\int_0^a f .
\]
故 $\int_{-1}^{1} \frac{t^3\cos t}{1 + t^4}\,\dd t = 0$，
无需求原函数（被积函数为奇函数），且
$\int_{-\pi}^{\pi} t^2\cos t\,\dd t = 2\int_0^\pi t^2\cos t\,
\dd t$。先检查对称性，再考虑各种技巧：
最快的积分，是那个从未被计算的积分。
\end{example}

\begin{example}[一眼识别导数]\label{ex:b1:integration:halfangle}
计算 $\displaystyle\int_0^{\pi/2} \frac{\dd x}{1 + \cos x}$。
利用半角公式 $1 + \cos x = 2\cos^2\frac x2$，将
被积函数化为 $\frac{1}{2}\bigl(1 + \tan^2\frac
x2\bigr)$，恰是 $\tan\frac x2$ 的导数：
\[
\int_0^{\pi/2} \frac{\dd x}{1 + \cos x}
= \Bigl[\tan\frac x2\Bigr]_0^{\pi/2} = \tan\frac\pi4 = 1 .
\]
无需任何换元机制，只需一种
把被积函数看作某个函数导数的习惯，其余
由微积分基本定理完成。（此类三角积分背后的系统工具，
即代换 $t = \tan\frac x2$，
属于标准工具箱，它建立在
\cref{thm:b1:integration:parts}~(2) 之上。）
\end{example}

\begin{example}[由积分定义的函数]\label{ex:b1:integration:definedby}
微积分基本定理能够构造函数。设
\[
F(x) = \int_0^x \eu^{-t^2}\,\dd t .
\]
没有任何初等函数的组合，其导数为
$\eu^{-t^2}$（刘维尔定理，在本阶段不加证明地承认）；但 $F$ 存在，
为 $C^1$ 类，且 $F'(x) = \eu^{-x^2} > 0$，严格递增，并为奇函数
（作代换 $t \mapsto -t$），且有界：当 $x \geq 1$ 时，
\[
F(x) - F(1) = \int_1^x \eu^{-t^2}\dd t
\leq \int_1^x \eu^{-t}\dd t \leq \eu^{-1} ,
\]
于是 $F \leq F(1) + \eu^{-1} \leq 1 + \eu^{-1}$。（极限的精确值
$\frac{\sqrt\pi}{2}$，借助二重积分计算，并在
本科三年级分册中完成。）变限积分的求导法则：$\frac{\dd}{\dd
x}\int_x^{x^2} \eu^{-t^2}\dd t = 2x\,\eu^{-x^4} - \eu^{-x^2}$.
末尾的要点：积分能\emph{构造}出新函数，
由旧函数出发，其全部性质都可由
被积函数读出——写不出的原函数，依然是
一个可被完全掌控的函数。
\end{example}

\begin{example}[不计算而估计]\label{ex:b1:integration:estimate}
积分 $R_n = \int_0^1 \frac{t^n}{1 + t}\,\dd t$ 没有
简洁的闭式，但由单调性可以精确地确定它们：在
$\intcc{0}{1}$ 上，$\frac12 \leq \frac{1}{1+t} \leq 1$，于是
\[
\frac{1}{2(n+1)} = \frac12\int_0^1 t^n\,\dd t
\;\leq\; R_n \;\leq\; \int_0^1 t^n \,\dd t = \frac{1}{n+1} :
\]
衰减的精确阶（$R_n \sim$ 某个 $\frac1n$ 的倍数，
事实上 $R_n \sim \frac{1}{2n}$）：两行推导，无需
求原函数。本章及下一章的周末问题，
恰好都基于这类框架——分析学家面对积分时，
首要的直觉应是\emph{估界}，而
只有此后，若有必要，才去计算它。
\end{example}

\begin{example}[平均值]\label{ex:b1:integration:average}
设连续函数 $f$ 在 $\intcc{a}{b}$ 上的\emph{平均值}为
$\frac{1}{b-a}\int_a^b f$。对正弦函数的一拱：
\[
\frac{1}{\pi}\int_0^\pi \sin t\,\dd t
= \frac{1}{\pi}\bigl[-\cos t\bigr]_0^\pi = \frac{2}{\pi}
\approx 0.637 :
\]
完整的正拱，其平均值并非 $\frac12$，而是
$\frac2\pi$——曲线在高处停留的时间
比三角形更长。由\cref{exo:b1:integration:11}（积分中值定理，
$g = 1$）可知平均值是一个\emph{值}：$\sin c =
\frac2\pi$，对某个 $c \in \intoo{0}{\pi}$ 成立。再由本章的黎曼
和，平均值是普通平均数的极限，
即 $n$ 个样本的平均，是连接数据离散平均与
信号连续平均的桥梁，这正是
积分进入物理学的方式。
\end{example}

\begin{theorem}[分部积分法；换元法]\label{thm:b1:integration:parts}
\begin{enumerate}
  \item 若 $u, v$ 为 $C^1$ 类，定义于 $\intcc{a}{b}$ 上：\index{分部
        积分法}
        \[
        \int_a^b u'v = \bigl[uv\bigr]_a^b - \int_a^b uv' .
        \]
  \item 若 $\varphi$ 为 $C^1$ 类，定义于 $\intcc{\alpha}{\beta}$ 上，且 $f$
        在 $\varphi(\intcc{\alpha}{\beta})$ 上连续：
        \index{换元法}
        \[
        \int_{\alpha}^{\beta} f\bigl(\varphi(t)\bigr)\,\varphi'(t)\,
        \dd t = \int_{\varphi(\alpha)}^{\varphi(\beta)} f(x)\, \dd x .
        \]
\end{enumerate}
\end{theorem}

\begin{proof}
(1) $(uv)' = u'v + uv'$；在 $\intcc{a}{b}$ 上积分，并对
$C^1$ 函数 $uv$ 应用微积分基本定理。

(2) 设 $F$ 为 $f$ 在像区间上的一个原函数
（见\cref{thm:b1:integration:ftc}）。于是 $(F \circ \varphi)' = (f
\circ \varphi)\,\varphi'$（链式法则），故两边都等于
$F(\varphi(\beta)) - F(\varphi(\alpha))$。
\end{proof}

\begin{example}\label{ex:b1:integration:parts}
$\int_0^1 t\,\eu^t \dd t = \bigl[t\,\eu^t\bigr]_0^1 - \int_0^1
\eu^t\dd t = \eu - (\eu - 1) = 1$。再作代换 $x =
\sin t$（$t \in \intcc{0}{\frac\pi2}$）：
\[
\int_0^1 \sqrt{1 - x^2}\, \dd x
= \int_0^{\pi/2} \cos t \cdot \cos t \, \dd t
= \int_0^{\pi/2} \frac{1 + \cos 2t}{2}\, \dd t = \frac\pi4 ,
\]
--- 恰为单位圆盘的四分之一，这正合几何直观。（线性化
方法在此起作用，见\cref{met:b1:complex:trig}）。
\end{example}

\begin{remark}[积分计算中的常见误区]
(i) \emph{代换在整个区间上须为 $C^1$ 函数}：
变换 $x = \frac1t$ 跨越 $0$ 时不合法；盲目地
用于 $\int_{-1}^{1}\frac{\dd x}{1 + x^2}$ 时，它“证明”该积分
等于它自身的相反数。当代换存在
奇点时，应先用沙勒关系将区间分段，
在每一段上分别代换，最后才合并。(ii) \emph{对数型
原函数需要绝对值}：$\int \frac{\dd x}{x - 2} =
\ln\abs{x - 2} + C$，在 $2$ 两侧分别成立 --- 写 $\ln(x - 2)$
（在 $\intoo{0}{1}$ 上）就是写一个
负数的对数；常数 $C$ 可在奇点的
两侧取不同的值。(iii) \emph{积分为零
并不意味着函数为零}：$\int_0^{2\pi}\sin = 0$。要断定 $f = 0$，
须先要求被积函数非负，
（见\cref{ex:b1:integration:chasleswork}）。(iv) \emph{黎曼和
须校准}：在 $\frac{b-a}{n}\sum f\bigl(a +
k\frac{b-a}{n}\bigr)$ 中，外部的步长与内部的取值点
必须对应同一分割 --- 常见错误是和式
$\sum_{k=1}^{n} f\bigl(\frac kn\bigr)$ \emph{不含}
因子 $\frac1n$，它发散而不是收敛到 $\int_0^1
f$。应用\cref{thm:b1:integration:riemann}前的检查清单：
提出因子 $\frac1n$，将求和项写成 $f$ 在 $\frac kn$ 处的值，
并指明 $f$，检验其连续性。
\end{remark}

\begin{example}[猜测、求导、调整]\label{ex:b1:integration:guessdiff}
求 $\int_1^x (\ln t)^2\,\dd t$？试猜一个
形如 $t\,P(\ln t)$ 的原函数，其中 $P$ 为多项式，并对其求导：
\[
\bigl(t\,P(\ln t)\bigr)' = P(\ln t) + P'(\ln t) .
\]
需要 $P(u) + P'(u) = u^2$，取 $P(u) = u^2 - 2u + 2$
（自 $u^2$ 起向下逐项比较系数）。于是
\[
\int_1^x (\ln t)^2\,\dd t
= \bigl[t\bigl((\ln t)^2 - 2\ln t + 2\bigr)\bigr]_1^x
= x(\ln x)^2 - 2x\ln x + 2x - 2 ,
\]
该结果亦可由两次分部积分得到。最后的要点在于：
对于形如
（$\ln t$ 的多项式）或（多项式乘 $\eu^{\lambda t}$）的被积函数，
其原函数具有相同形式 --- 对一个设定形式的猜测求导，
可将积分化为关于系数的线性代数问题，
比反复分部积分更快，也更不易出错。
\end{example}

\begin{example}[回旋镖积分]\label{ex:b1:integration:boomerang}
计算 $I = \int_0^{\pi/2} \eu^x \cos x\,\dd x$。
两次分部积分，每次都对三角因子求导：
\[
I = \bigl[\eu^x \sin x\bigr]_0^{\pi/2}
- \int_0^{\pi/2} \eu^x \sin x\,\dd x
= \eu^{\pi/2} - J,
\qquad
J = \bigl[-\eu^x\cos x\bigr]_0^{\pi/2}
+ \int_0^{\pi/2} \eu^x\cos x\,\dd x = 1 + I .
\]
积分回到了自身：$I = \eu^{\pi/2} - 1 - I$，
由此
\[
I = \frac{\eu^{\pi/2} - 1}{2} .
\]
末了的要点：当被积函数是两个求导后
仍能复现自身的函数之积时
（$\eu^{ax}$、$\cos bx$、$\sin bx$），两次分部积分
会得到关于未知积分的线性方程 --- 解此方程，
而不是继续积分；等价地，经由
$\eu^{(a + \iu b)x}$（见\cref{ch:b1:complex}）并取实部。
两条路径给出相同答案，验证二者相符
是一项免费的合理性检验。
\end{example}

\section{黎曼和}

\begin{theorem}[黎曼和]\label{thm:b1:integration:riemann}
设 $f$ 在 $\intcc{a}{b}$ 上连续，则\index{黎曼和}
\[
S_n = \frac{b - a}{n} \sum_{k=0}^{n-1}
f\Bigl(a + k\,\frac{b-a}{n}\Bigr)
\xrightarrow[n \to \infty]{} \int_a^b f(t)\, \dd t ,
\]
对子区间内任意取值点，结论同样成立。
\end{theorem}

\begin{proof}
$S_n$ 是阶梯函数 $\varphi_n$ 的积分，其值为
$f(a + k\frac{b-a}{n})$，即在第 $k$ 个子区间上。
任给 $\varepsilon > 0$，由一致连续性（海涅）可得 $\delta$；当
$n > \frac{b-a}{\delta}$ 时，每个子区间中的点与其取值点的距离都
不超过 $\delta$，故 $\abs{f - \varphi_n} \leq
\varepsilon$ 在 $\intcc{a}{b}$ 上成立，于是
\[
\Bigl| \int_a^b f - S_n \Bigr|
= \Bigl| \int_a^b (f - \varphi_n) \Bigr|
\leq (b-a)\,\varepsilon . \qedhere
\]
\end{proof}

\begin{omfigure}
\begin{tikzpicture}
\begin{axis}[omaxis, width=8.6cm, height=5.4cm,
    xmin=0, xmax=1.15, ymin=0, ymax=1.7,
    xtick={0,1}, ytick=\empty]
  \draw[fill=omDef!15, draw=omDef!60, thin]
    (axis cs:0,0)      rectangle (axis cs:0.125,0.4)
    (axis cs:0.125,0)  rectangle (axis cs:0.25,0.4156)
    (axis cs:0.25,0)   rectangle (axis cs:0.375,0.4625)
    (axis cs:0.375,0)  rectangle (axis cs:0.5,0.5406)
    (axis cs:0.5,0)    rectangle (axis cs:0.625,0.65)
    (axis cs:0.625,0)  rectangle (axis cs:0.75,0.7906)
    (axis cs:0.75,0)   rectangle (axis cs:0.875,0.9625)
    (axis cs:0.875,0)  rectangle (axis cs:1,1.1656);
  \addplot[omThm, very thick, domain=0:1.05, samples=60]
    {0.4 + x^2};
\end{axis}
\end{tikzpicture}

{\small 取 $n = 8$ 个矩形的左黎曼和：分划
逐步加细时，一致连续性迫使阶梯形面积趋向
$\int_a^b f$。}
\end{omfigure}

\begin{example}\label{ex:b1:integration:riemannexample}
$\displaystyle\sum_{k=1}^{n} \frac{1}{n + k} = \frac 1n
\sum_{k=1}^{n} \frac{1}{1 + k/n} \xrightarrow[n\to\infty]{}
\int_0^1 \frac{\dd x}{1 + x} = \ln 2$：一个初等估计无法触及的极限，
但写成黎曼和则一目了然。
\end{example}

\begin{example}[第二个黎曼和，附配凑]\label{ex:b1:integration:riemann2}
求 $\lim_{n\to\infty} \sum_{k=1}^{n} \dfrac{n}{(n+k)^2}$。
配凑：
\[
\sum_{k=1}^{n} \frac{n}{(n+k)^2}
= \frac{1}{n}\sum_{k=1}^{n} \frac{n^2}{(n+k)^2}
= \frac1n \sum_{k=1}^{n} \frac{1}{\bigl(1 + \frac kn\bigr)^2} ,
\]
即连续函数 $f(x) = \frac{1}{(1+x)^2}$ 在
$\intcc{0}{1}$ 上的黎曼和，于是极限为
\[
\int_0^1 \frac{\dd x}{(1 + x)^2}
= \Bigl[-\frac{1}{1+x}\Bigr]_0^1 = \frac12 .
\]
小结：全部技巧在于中间一行——
将求和项化为 $f(\frac kn)$ 的形式，代价是
恰好提出一个因子 $\frac1n$；形式一旦正确，
定理完成分析，微积分基本定理
则完成计算。
\end{example}

\begin{remark}[积分的后续应用]
本章的每一种构造都有后续。黎曼和
将在\cref{ch:b1:series}中再次出现，作为级数与积分之间的桥梁
（比较 $\sum \frac{1}{n^\alpha}$ 与 $\int
\frac{\dd t}{t^\alpha}$）；积分余项给出
\cref{ch:b1:taylor}中泰勒公式最精确的形式；
基于 $\sup$ 的定义，是本科三年级分册中勒贝格积分的原型，
其中同样的三条性质
（线性、单调性、收敛定理）在
一类远为广泛的函数上重新建立。下面的周末问题
把分部积分化为算术，即证明
$\pi^2$ 的无理性。
\end{remark}

\section{练习}

\begin{exercise}[$\star$]\label{exo:b1:integration:1}
计算：$\displaystyle\int_0^1 \frac{\dd x}{x^2 - 4}$
\emph{（部分分式，\cref{ch:b1:fractions}）}；
$\displaystyle\int_1^{\eu} \ln t \,\dd t$；
$\displaystyle\int_0^{\pi} t \sin t\, \dd t$。
\end{exercise}

\begin{exercise}[$\star$]\label{exo:b1:integration:2}
计算 $\displaystyle\int_0^{1} \frac{t}{(t^2+1)^2}\,\dd t$
（换元法）和
$\displaystyle\int_0^{\pi/2} \cos^3 t\, \dd t$（可写成 $\cos^3 =
\cos(1 - \sin^2)$）。
\end{exercise}

\begin{exercise}[$\star$]\label{exo:b1:integration:3}
用黎曼和求下列极限：
\[
u_n = \sum_{k=1}^{n} \frac{n}{n^2 + k^2},
\qquad
v_n = \frac{1}{n}\sqrt[n]{\frac{(2n)!}{n!\,n^n}}
\ \emph{（取对数）}。
\]
\end{exercise}

\begin{exercise}[$\star$]\label{exo:b1:integration:4}
设 $f$ 在 $\intcc{0}{1}$ 上连续。计算 $\lim_{n\to\infty}
\int_0^1 x^n f(x)\,\dd x$。\emph{（将 $\intcc{0}{1}$ 截断于 $1 -
\delta$ 处。）}
\end{exercise}

\begin{exercise}[$\star\star$]\label{exo:b1:integration:5}
（柯西–施瓦茨不等式）设 $f, g$ 在 $\intcc{a}{b}$ 上连续，证明
\[
\Bigl(\int_a^b fg\Bigr)^{\!2} \leq \int_a^b f^2 \cdot \int_a^b g^2 ,
\]
将 $\int_a^b (f + \lambda g)^2 \geq 0$ 展开为
关于 $\lambda$ 的二次式。何时等号成立？
\end{exercise}

\begin{exercise}[$\star\star$]\label{exo:b1:integration:6}
设 $f$ 为 $\R$ 上的连续函数，且以 $T$ 为周期。证明 $\int_a^{a+T}
f$ 与 $a$ 无关，并证明 $\frac1x \int_0^x f(t)\,\dd t
\to \frac 1T \int_0^T f$，当 $x \to +\infty$ 时成立。
\end{exercise}

\begin{exercise}[$\star\star$]\label{exo:b1:integration:7}
设 $f$ 在 $\intcc{0}{1}$ 上连续，且 $\int_0^1 f = \frac12$，
证明 $f$ 在 $\intcc{0}{1}$ 中存在不动点。\emph{（对
$f(x) - x$ 积分，并利用严格正性，
借助 \cref{thm:b1:integration:props}~(4) 的逆否命题，
并结合介值定理。）}
\end{exercise}

\begin{exercise}[$\star\star\star$]\label{exo:b1:integration:8}
（沃利斯积分）设 $W_n = \int_0^{\pi/2} \sin^n t\,\dd t$。
\begin{enumerate}
  \item 用分部积分证明递推式 $n W_n = (n-1) W_{n-2}$（$n \geq 2$），
        并求 $W_0, W_1$，再求 $W_{2p}$ 与 $W_{2p+1}$
        的封闭形式。
  \item 证明 $(W_n)$ 单调递减，且 $\frac{W_{n+1}}{W_n}
        \to 1$；证明量 $(n+1)\,W_{n+1} W_n$ 为常数，
        其值为 $\frac\pi2$，并由此推出等价式
        $W_n \sim \sqrt{\dfrac{\pi}{2n}}$。
\end{enumerate}
\end{exercise}

\begin{exercise}[$\star\star\star$]\label{exo:b1:integration:9}
（尼文：$\pi$ 为无理数）设 $\pi = \frac ab$，其中 $a, b \in
\N^*$，对待定的 $n$，令
\[
P(x) = \frac{x^n (a - bx)^n}{n!},
\qquad
I_n = \int_0^{\pi} P(x)\sin x\, \dd x .
\]
\begin{enumerate}
  \item 证明 $0 < I_n \leq \pi\,\frac{(\pi a)^n}{n!}$，
        该上界 $< 1$，只要 $n$ 充分大。
  \item 证明 $P$ 及其所有各阶导数
        在 $0$ 与 $\pi = \frac ab$ 处均取\emph{整数}值。\emph{（二项式
        展开：$P$ 的系数乘以 $k!$ 均为整数，
        当 $k \geq n$ 时；并有 $P(\pi - x) = P(x)$。）}
  \item 令 $Q = P - P'' + P^{(4)} - \dots$（有限和），验证
        $\bigl(Q'\sin x - Q\cos x\bigr)' = P \sin x$，并由此推出
        $I_n = Q(\pi) + Q(0)$ 为整数。
  \item 得出结论。
\end{enumerate}
\end{exercise}

\begin{exercise}[$\star\star\star$]\label{exo:b1:integration:10}
设 $f$ 为 $C^1$ 函数，定义于 $\intcc{a}{b}$ 上，证明下述黎曼--勒贝格型
极限
\[
\int_a^b f(t)\sin(\lambda t)\,\dd t \xrightarrow[\lambda \to
+\infty]{} 0
\]
用分部积分证明此极限。随后对仅为
连续的 $f$，利用阶梯函数的一致逼近再证明之
（见\cref{thm:b1:integration:approx}）。
\end{exercise}

\begin{exercise}[$\star\star$]\label{exo:b1:integration:11}
（积分中值定理）设 $f, g$ 在
$\intcc{a}{b}$ 上连续，且 $g \geq 0$。证明存在 $c \in
\intcc{a}{b}$，使得
\[
\int_a^b f(t)\,g(t)\,\dd t = f(c)\int_a^b g(t)\,\dd t ,
\]
并通过一个例子说明条件 $g \geq 0$ 不可
省去。
\end{exercise}

\begin{exercise}[$\star\star\star$]\label{exo:b1:integration:12}
（矩迫使出现零点）设 $f$ 在 $\intcc{a}{b}$ 上连续，
且满足
\[
\int_a^b f(t)\,t^k\,\dd t = 0 \qquad \text{对 } k = 0, 1,
\dots, n .
\]
证明 $f$ 在 $n + 1$ 个互不相同的点处为零，
这些点位于 $\intoo{a}{b}$ 内。\emph{（若 $f$ 仅在 $z_1 < \dots
< z_m$ 处变号，且 $m \leq n$，对 $f$ 与 $P(t) = (t - z_1)
\cdots (t - z_m)$ 之积作积分，并利用严格正性。）}
\end{exercise}

\begin{remark}[本分册内的展望]
后面三章直接依赖本章。
\Cref{ch:b1:taylor} 给出积分型余项 --- 三个泰勒公式中最精细的一个，
恰为 $n$ 次迭代的分部积分。
\Cref{ch:b1:series} 将以积分夹逼和式的框架，
转化为 $\sum n^{-\alpha}$ 的决定性判别法，
其周末问题将该框架进一步细化为欧拉
常数。\Cref{ch:b1:curves} 使积分具有几何意义：
参数曲线的弧长为 $\int \sqrt{x'(t)^2 +
y'(t)^2}\,\dd t$，是连续函数在
闭区间上的积分 --- 正是本章构造的对象，无需
反常积分理论。最常被复用的事实将是最朴素的：
$\bigl|\int f\bigr| \leq (b - a)\sup\abs{f}$，此不等式
能将每个逐点估计化为积分估计。
\end{remark}

\section{问题：积分无理性机器}

\begin{problem}[{周末问题 --- $\eu$ 与 $\pi^2$ 均为无理数，
$\eu$ 精确到小数点后六位，且 $\frac{22}{7} > \pi$，附
证明}]
\label{pb:b1:integration:1}
全题只靠一个机制驱动：\emph{一个式子，它必须是正整数，
却可证明小于 $1$，
这样的式子不可能存在}。\Cref{exo:b1:integration:9}（尼文）曾用它一次证明
$\pi \notin \Q$，此处将其规模化。此“机器”
需要三个部件：一个\emph{整数性}输入（精心选取的多项式的端点值），
一个\emph{小量性}输入（一个因子
$\frac{1}{n!}$ 使积分极小），以及一座\emph{桥梁}
（分部积分）联结二者。证明 $\eu$ 为无理数，
并以有保证的误差算出其值，证明勒让德的
更强的定理，即 $\pi^2$ 为无理数，并以
分析学中最迷人的积分作结：$\int_0^1 \frac{x^4(1 -
x)^4}{1 + x^2}\dd x = \frac{22}{7} - \pi$，它
以手算夹逼 $\pi$。\index{无理性！
$\pi^2$}\index{无理性！$\eu$}

\medskip
\noindent\textbf{第一部分 --- 燃料。}
\begin{enumerate}
  \item 证明 $\dfrac{c^{\,n}}{n!} \to 0$，对每个固定的 $c
        > 0$ \emph{（当 $n \geq 2c$ 之后，每一步至少使该项减半
        ）}。
  \item （贝塔积分）对 $m$ 用归纳法，并结合
        分部积分证明：
        \[
        \int_0^1 x^{\,k}\,(1 - x)^{\,m}\,\dd x =
        \frac{k!\,m!}{(k + m + 1)!}
        \qquad (k, m \in \N).
        \]
  \item 由此推出 $\displaystyle\int_0^1 \bigl(x(1-x)\bigr)^n \dd x
        = \frac{1}{(2n+1)\binom{2n}{n}}$，并 --- 与
        界 $x(1 - x) \leq \frac14$ 比较 --- 由此得
        估计 $\binom{2n}{n} \geq \dfrac{4^n}{2n+1}$，与
        \cref{pb:b1:seq:1}中的 $\binom{2n}{n}^{1/n} \to 4$ 相吻合。
  \item 证明下文两次用到的小量引理：对每个
        连续函数 $g > 0$，定义于 $\intoo{0}{1}$ 上，
        \[
        0 < \int_0^1 \bigl(x(1-x)\bigr)^n g(x)\,\dd x
        \leq \frac{\sup_{\intcc{0}{1}}\abs g}{4^{\,n}} .
        \]
\end{enumerate}

\medskip
\noindent\textbf{第二部分 --- $\eu$：无理性，再求六位
小数。} 令 $A_n = \displaystyle\int_0^1 x^n \eu^x \dd x$。
\begin{enumerate}[resume]
  \item 计算 $A_0$ 与 $A_1$，证明递推关系 $A_n = \eu
        - n A_{n-1}$，并证明界 $0 < A_n \leq
        \dfrac{\eu}{n+1}$。
  \item 用归纳法证明 $A_n = \alpha_n + \beta_n \eu$，
        其中 $\alpha_n, \beta_n \in \Z$。
  \item 由此推出 $\eu$ 为无理数 \emph{（若 $\eu = \frac
        pq$，则 $q A_n$ 为一个落在
        $\intoo{0}{1}$ 内的整数，对充分大的 $n$ 成立）}。比较
        \cref{exo:b1:seq:9}的证明：结论相同，所用“燃料”不同。
  \item 用归纳法并结合分部积分证明精确的
        余项公式
        \[
        \eu = \sum_{k=0}^{n} \frac{1}{k!} + R_n,
        \qquad
        R_n = \frac{1}{n!}\int_0^1 (1 - t)^n\,\eu^{\,t}\,\dd t,
        \qquad
        \frac{1}{(n+1)!} \leq R_n \leq \frac{\eu}{(n+1)!} .
        \]
  \item 取 $n = 9$，估计 $R_9$ 时利用 $\eu < 2.75$（由
        \cref{ex:b1:seq:e}中的 $b_2 = 2.75$ 得出），计算和式，
        并得出有保证的夹逼 $2.7182818 \leq
        \eu \leq 2.7182823$ --- 六位小数，$\eu \approx
        2.718282$，并给出证明。
\end{enumerate}

\medskip
\noindent\textbf{第三部分 --- 勒让德定理：$\pi^2$ 为
无理数。} 设 $f(x) = \dfrac{x^n (1 - x)^n}{n!}$，并假设
$\pi^2 = \frac ab$，其中 $a, b \in \N^*$。
\begin{enumerate}[resume]
  \item 证明 $f(1 - x) = f(x)$，且 $0 < f \leq
        \dfrac{1}{4^n\,n!}$ 在 $\intoo{0}{1}$ 上成立。
  \item 证明 $f^{(k)}(0)$ 与 $f^{(k)}(1)$ 对每个
        $k \geq 0$ 都是整数 \emph{（将 $x^n(1-x)^n$ 按整
        系数展开；$\frac{k!}{n!} \in \Z$ 对 $k \geq n$ 成立；
        再利用对称性）}。
  \item 令
        \[
        G = b^{\,n} \sum_{k=0}^{n} (-1)^k\,
        \pi^{2n - 2k} f^{(2k)} .
        \]
        证明 $G(0)$ 与 $G(1)$ 为整数 \emph{（每一项
        $b^n \pi^{2n-2k} = a^{\,n-k}\,b^{\,k}$）}。
  \item 验证裂项相消恒等式 $G'' + \pi^2 G = b^n \pi^{2n+2} f
        = \pi^2 a^n f$，于是
        \[
        \frac{\dd}{\dd x}\Bigl(G'(x)\sin \pi x - \pi\,G(x)\cos
        \pi x\Bigr) = \pi^2 a^n f(x)\sin \pi x .
        \]
  \item 在 $\intcc{0}{1}$ 上积分，并得出
        \[
        \pi\,a^n \int_0^1 f(x)\sin(\pi x)\,\dd x = G(0) + G(1)
        \in \Z ,
        \]
        一个\emph{正}整数，且其值不超过 $\dfrac{\pi
        a^n}{4^n\,n!}$。
  \item 结合第 1 题，证明 $\pi^2$ 为无理数
        （勒让德，1794 年），并说明这加强了
        \cref{exo:b1:integration:9}的结论：为什么
        $\pi^2$ 的无理性蕴含 $\pi$ 的无理性，而反之不然？
\end{enumerate}

\medskip
\noindent\textbf{第四部分 --- 理解这台“机器”。}
\begin{enumerate}[resume]
  \item 在第二、三部分中找出两股对立的力量（端点数据的整数性；积分的解析小量性）
        以及连接二者的“桥梁”。然后解释，为何因子
        $\frac{1}{n!}$ 在 $f$ 中起决定性作用：若去掉它，
        整数性依然保留，
        但哪一个不等式失效，以及
        对哪些所声称的分数 $\frac ab$，
        证明随之失败？
  \item 有效性：设有人声称 $\pi^2 = \frac ab$
        其中 $a \leq 10$。证明矛盾已在
        $n = 7$ 时出现：计算 $\pi\,(10/4)^7/7! \approx
        0.38 < 1$。机器不只是反驳，而是
        借助一个固定、可计算的阶段进行反驳。
  \item 无条件地检验“桥梁”：用两次分部积分证明
        下式成立：
        \[
        \int_0^1 x(1 - x)\sin(\pi x)\,\dd x = \frac{4}{\pi^3},
        \]
        并与第 14 题在 $n = 1$ 时的结果相互印证（保持 $\pi^2$ 为符号形式；
        裂项相消恒等式为 $\pi^3
        \int_0^1 f_1 \sin \pi x = -(f_1''(0) + f_1''(1)) = 4$）。
  \item 是什么使 $\eu^x$ 与 $\sin \pi x$ 有资格充当
        “机器”的核？请指出其性质（二者各自满足
        一个常系数线性微分方程，
        因此反复分部积分会循环回到起点），并指出前沿：同一台
        机器经赫尔米特与林德曼改进，可证明
        $\eu$ 与 $\pi$ 为\emph{超越数}
        --- 这已超出本册的范围。
\end{enumerate}

\medskip
\noindent\textbf{第五部分 --- $\frac{22}{7}$ 与 $\pi$ 的比较与
寓意。}
\begin{enumerate}[resume]
  \item 验证下列多项式除法等式
        \[
        \frac{x^4(1-x)^4}{1 + x^2}
        = x^6 - 4x^5 + 5x^4 - 4x^2 + 4 - \frac{4}{1 + x^2},
        \]
        并由此推出著名的恒等式
        \[
        \int_0^1 \frac{x^4 (1-x)^4}{1 + x^2}\,\dd x
        = \frac{22}{7} - \pi .
        \]
  \item 被积函数为正，由此得 $\pi <
        \frac{22}{7}$。再将 $\frac{1}{1+x^2}$ 放缩到
        $\frac12$ 与 $1$ 之间，并利用 $\int_0^1 (x(1-x))^4 =
        \frac{1}{630}$（第 3 题），证明
        \[
        \frac{22}{7} - \frac{1}{630} \;\leq\; \pi \;\leq\;
        \frac{22}{7} - \frac{1}{1260} ,
        \]
        即 $3.14126 \leq \pi \leq 3.14207$：
        两位正确的小数，均可手算得出。
  \item 推广：将 $x^{4m}(1-x)^{4m}$ 除以 $1 + x^2$，
        证明余式为常数 $(-4)^m$ \emph{（在模 $x^2 + 1$ 意义下计算
        ：$(1-x)^2 \equiv -2x$）}，由此推出
        有理数 $r_m$，满足
        \[
        \abs{\pi - r_m} \leq 4^{\,1 - 5m} ,
        \]
        并验证 $m = 1$ 时复现第 20--21 题的结果。
  \item 将这些有理数与
        \cref{pb:b1:derivative:1}中的逼近理论相对照：计算 $\abs{\pi -
        \frac{22}{7}} \approx 1.26\cdot10^{-3}$，与狄利克雷
        保证值 $\frac{1}{49}$ 相比较，并引用
        $\abs{\pi - \frac{355}{113}} \approx 2.7\cdot10^{-7}$
        与 $\frac{1}{113^2} \approx 7.8\cdot10^{-5}$ 相对照：
        对 $\pi$ 存在异常优良的有理逼近
        --- 这与之相容，因为尚不知 $\pi$ 是否
        为坏逼近数。
  \item （抽象化的“整数陷阱”）证明统一以上全部内容的引理：
        \emph{若 $x \in \R$，且存在
        整数 $a_n, b_n$ 满足 $0 < \abs{a_n + b_n x} \to 0$，
        则 $x$ 为无理数。}列出其在本题
        中、在\cref{exo:b1:integration:9}中、
        在\cref{exo:b1:seq:9}中，以及在\cref{pb:b1:derivative:1}中的实例。
  \item 综合归纳，每项各用一句话：(i) 该“机器”的三个
        部分及其在本章工具箱中的
        位置；(ii) 积分所提供而
        \cref{pb:b1:derivative:1}中的中值定理所不能提供的；
        (iii) 所提取结果的清单（两个
        无理性结论、一个精确到六位小数的常数、一个
        对 $\pi$ 的夹逼、一个二项式界）；(iv) 前沿
        （赫尔米特、林德曼；以及同一陷阱在
        $\zeta(2)$ 与 $\zeta(3)$ 上运行，
        二十世纪算术中的情形）。
\end{enumerate}
\end{problem}
